\documentclass[10pt,a4paper,twocolumn]{article}
\usepackage[a4paper,margin=1.6cm,columnsep=0.6cm]{geometry}
\usepackage{amsmath}
\usepackage{txfonts}
\usepackage[T1]{fontenc}
\usepackage[protrusion=false]{microtype}
\usepackage{booktabs}
\usepackage{enumitem}
\usepackage{titlesec}
\setlist{nosep,leftmargin=1.1em}
\titleformat{\section}{\normalsize\bfseries}{\thesection}{0.5em}{}
\titlespacing*{\section}{0pt}{5pt}{2pt}
\setlength{\parskip}{2pt}
\setlength{\parindent}{0pt}
\pagestyle{empty}

\begin{document}
\twocolumn[{\centering
{\large\bfseries When Does Graph Expansion Help? Issue-Adaptive Retrieval Routing for Repository-Level Code Localization}\\[3pt]
{\normalsize Midway report, 31 October 2026}\\[2pt]
{\small Naman Singhal (2024114013) \quad Yash More (2024114004) \quad Aaditya Bhatiya (2023114012)}\\[6pt]}]

\section{Where we are}
Stages 1 to 4 of the proposal's timeline are done, about a week ahead of
plan. Phase A (the measurement study) is complete, as is the first
version of the rule-based router. What remains is testing the router on
held-out repositories, followed by the figures and the write-up.

\section{What is built}
\begin{itemize}
\item Data: the Python subset of SWE-Explore is 500 issues (410 from
SWE-bench Verified and 90 from SWE-bench Pro), joined with SWE-bench for the
issue text and pre-fix commit.
\item Pipeline: pre-fix snapshots, function, class and module chunks, a BM25
seed and a typed repository graph (with call, import, containment and
inheritance edges) adapted from LocAgent's builder. Our version is about 5 times
faster and gave exactly the same edges on the five repositories we checked.
\item Scorer: SWE-Explore's released scorer, vendored unchanged, alongside our
primary metrics (below). Both pass hand-computed golden examples.
\item Actions, all cut to the same line budget: stop (seed only), dense
(frozen CodeRankEmbed), BM25 and dense fusion (RRF, $k=100$ tuned on dev) and
one-hop expansion along calls, imports or inheritance.
\item Phase A: all 6 actions on all 500 issues at budgets of 100, 300 and
500 lines, plus 16 cheap post-seed features, the signature table and a
Tier 1 rule procedure fixed before any held-out scoring.
\end{itemize}

\section{Phase A results (nDCG at 300 lines, 500 issues)}
\begin{center}\small
\begin{tabular}{lccc}
\toprule
Action & nDCG & helps & hurts \\
\midrule
stop (BM25 seed) & .172 & & \\
expand(inherit) & .172 & 3\% & 2\% \\
expand(imports) & .170 & 38\% & 42\% \\
expand(calls) & .178 & 41\% & 37\% \\
dense & .152 & 45\% & 48\% \\
fusion & .184 & 52\% & 35\% \\
per-issue oracle & .258 & & \\
\bottomrule
\end{tabular}\end{center}
``Helps'' and ``hurts'' count issues where the action beats or loses to the
seed. At 300 lines only fusion reliably beats the seed on average and it
raises SWE-bench file recall from .628 to .706. Graph expansion barely moves
the mean, although at 100 lines call expansion is the best fixed action. Per issue, though, expansion helps about 40\% of issues and hurts
about 40\%. Choosing the best action per issue would be 40\% above the best
fixed action (fusion). The oracle's choices are spread evenly across five
actions, which supports the premise that the right action depends on the
issue.

Of the proposal's five expectations, only the hub one is supported. Even that
holds only as a contrast: call expansion gains .017 [.009, .031] when the top hit is
weakly connected but $-$.014 [$-$.036, .001] when it is a hub. In the other four, a
traceback did not make call expansion help, naming the target file did not
make expansion useless, seed concentration predicted nothing and
disagreement between BM25 and dense did not favour fusion. The strongest
signal is the graph degree of the top seed hit. One rule recurs in all 8
folds: when the top hit has many call edges and dense disagrees with BM25,
keep the seed and expand along calls, since fusion does worst on these
issues (stop .202, expand(calls) .213, fusion .161 on 64 issues).

\section{Changes from the proposal}
\begin{itemize}
\item The Python subset has only 15 repositories and django holds 41\% of
the issues. We use a
fixed dev group (xarray, pytest, openlibrary: 56 issues) for all tuning and
leave-one-repository-out folds (8 folds) for testing.
\item The ground truth is often whole files read by agents (median 1,160
lines), so even a perfect ranker averages only 39.5\% recall at 300 lines.
nDCG is the primary metric. We follow the SWE-Explore paper's definition,
in which the ideal list is built from the ground truth. The released code
instead rescores the returned list, so a seed with one hit at rank 1 gets
1.0. We report both.
\item Ranges past the end of a file are clipped in our metrics. The official
scorer keeps them in the denominator.
\item Edge direction (outgoing only, or both ways) is tuned on dev with the
decay. Both ways is best for calls and imports. Inheritance edges are rare,
so expand(inherit) is almost identical to stop.
\item The graph builder uses Python's \texttt{ast} module. The proposal said
tree-sitter.
\item Related work now includes the closest precedent we found, from classic
IR: selective pseudo-relevance feedback and query performance prediction
(Amati et al.\ 2004, Cronen-Townsend et al.\ 2004), which also predict from
first-pass results whether expansion will help.
\end{itemize}

\section{Next steps}
Stage 5 (2 to 8 November): fit the rule procedure on each fold's training
repositories, route the held-out repository and report GapClosed against
the best fixed action and the oracle, with repository-clustered intervals and
cost. The router with up to 5 rules is the main result. The same router cut
after its first rule is a secondary analysis, declared now. Tier 2 (a small
classifier) only if Stage 5 finishes early.

\end{document}
